% Partie : sets-functions-relations
% Chapitre : size-of-sets
% Section : introduction

\documentclass[../../../include/open-logic-section]{subfiles}

\begin{document}

\olfileid[fr]{sfr}{siz}{int}

\olsection{Introduction}

Lorsque Georg Cantor développa la théorie des ensembles dans les années
1870, il cherchait notamment à faire admettre l’idée d’une collection
infinie : un infini en acte, comme l’auraient dit les penseurs médiévaux.
Son étude de la \emph{taille} des ensembles joua un rôle essentiel
dans cette démarche. Si $a$, $b$ et $c$ sont tous distincts, l’ensemble
$\{a, b, c\}$ est intuitivement \emph{plus grand} que $\{a, b\}$.
Mais qu’en est-il des ensembles infinis ? Ont-ils tous la même taille ?
Il se trouve que non.

La première idée importante est celle d’énumération. On peut dresser
la liste des !!{element}s de tout ensemble fini. Pour certains ensembles
infinis, on peut aussi dresser une telle liste, à condition de permettre
qu’elle soit elle-même infinie. Ces ensembles sont dits !!{enumerable}s.
Le résultat surprenant de Cantor, que nous comprendrons pleinement
à la fin de ce chapitre, est que certains ensembles infinis ne sont pas
!!{enumerable}s.

\end{document}
